%% Figure: overview of \tool (Protocol B). Steps left to right; the dashed
%% back-edge is the outer loop over pinned OpenSSL snapshots in run_scan.sh,
%% not an iteration inside one file's analysis.
\begin{tikzpicture}[
    x=1cm, y=1cm, font=\scriptsize,
    box/.style={draw, rounded corners=2pt, align=center, inner sep=2pt,
                minimum height=10mm, text width=17mm, fill=white},
    outbox/.style={draw, align=left, inner sep=2pt, fill=black!6,
                text width=22mm, font=\tiny},
    flow/.style={-{Stealth[length=1.6mm]}, line width=0.5pt},
    back/.style={-{Stealth[length=1.6mm]}, line width=0.5pt, dash pattern=on 2pt off 1.5pt},
  ]
  \node[box, text width=15mm] (in) at (-0.3,0) {C source\\+ rule (JSON)};
  \node[box] (s1) at (2.1,0) {\textbf{1} parse \&\\ build CFG};
  \node[box] (s2) at (4.3,0) {\textbf{2} find release\\ sites; search\\ paths to exit};
  \node[box] (s3) at (6.5,0) {\textbf{3} classify \&\\ extract witness};
  \node[outbox] (o) at (4.3,-1.8) {per-site result:\\ conforming /\\ violation /\\ suppressed\\ + witness path};
  \draw[flow] (in) -- (s1);
  \draw[flow] (s1) -- (s2);
  \draw[flow] (s2) -- (s3);
  \draw[flow] (s3.south) |- (o.east);
  \draw[back] (s3.north) -- ++(0,0.5) -| node[pos=0.3, above] {next snapshot} (s1.north);
\end{tikzpicture}
